\Needspace{20\baselineskip}
\fancyhead[L]{\small KOMET VON 1807.}
\begin{center}
KOMET VON 1807.
\end{center}

\begin{center}
\fontsize{10}{12}\selectfont\begin{tabular}{c}
\hline
Göttingische gelehrte Anzeigen. Stück 53. Seite 521..523. 2. April 1808. \\
\hline
\end{tabular}
\end{center}

Einem Brief an Prof. \textsc{Gauss} vom 18. März zufolge hatte Inspektor \textsc{Bessel} in Lilienthal das Glück, den letzten Kometen noch bis gegen Ende Februar verfolgen zu können. Wir beeilen uns, diese wertvollen Beobachtungen hier mitzuteilen.


\begin{center}\fontsize{10}{12}\selectfont\setlength{\tabcolsep}{3pt}\renewcommand{\arraystretch}{1.15}\begin{tabular}{lrrr}
1808 & Mittlere Zeit & \shortstack{Scheinbare\\Rektaszension} & \shortstack{Scheinbare\\nördl. Deklination}\\\hline
Jan. 4 & $9^{\text{u}}\,49^{\text{m}}\,40^{\text{s}}$ & $325^\circ\;17'\;16''1$ & $47^\circ\;13'\;11''1$\\
4 & $11\;34\;32$ & $325\;21\;47.6$ & $47\;13\;40.9$\\
12 & $6\;58\;59$ & $333\;3\;58.2$ & $47\;40\;58.8$\\
21 & $7\;26\;15$ & $341\;19\;13.0$ & $47\;59\;15.3$\\
23 & $6\;49\;6$ & $343\;2\;37.9$ & $48\;2\;5.2$\\
Febr. 19 & $7\;20\;6$ & $3\;43\;35.7$ & $48\;21\;17.1$\\
20 & $7\;23\;15$ & $4\;24\;3.7$ & $48\;21\;33.5$\\
24 & $7\;35\;16$ & $7\;3\;11.1$ & $48\;23\;38.3$\\
\end{tabular}\end{center}

\par\indent Nach Herrn \textsc{Bessel}s eigenem Urteil verdient unter den im Februar beobachteten Deklinationen allein die vom 20. Vertrauen.

W\SourceLeaf{306}ir fügen diesen Beobachtungen auch die späteste von Dr. \textsc{Olbers} in Bremen hinzu. Die früheren Beobachtungen dieses Astronomen sind bereits aus der Monatl. Correspondenz bekannt:

\begin{quote}
Febr. 14. $7^{\text{u}}\,30^{\text{m}}\,36^{\text{s}}$ mittl. Z.\quad Rektasz. $0^\circ\;15'\;53''$,\quad nördl. Deklination $48^\circ\;18'\;10''$
\end{quote}

\par\indent Herr \textsc{Bessel} hat diese letzten Beobachtungen benutzt, um seine parabolischen Elemente noch einmal zu verbessern. Dabei erhielt er die folgenden Ergebnisse, die sich nur wenig von den Elementen des Herrn Prof. \textsc{Gauss} unterscheiden:


\begin{center}\fontsize{9}{11}\selectfont\setlength{\tabcolsep}{3pt}\renewcommand{\arraystretch}{1.15}\begin{tabular}{lr}
\shortstack[l]{Durchgang durch das Perihel,\\Pariser Zeit 1807 September} .\;.\;.\; & $18{,}82718$\\
Neigung der Bahn .\;.\;.\;.\;.\;.\;.\;.\;.\;.\;.\;.\;.\;.\;.\;.\;.\; & $63^\circ\;14'\;28''1$\\
Aufsteigender Knoten .\;.\;.\;.\;.\;.\;.\;.\;.\;.\;.\;.\;.\;.\;.\;.\; & $266\;36\;51{,}7$\\
Perihel .\;.\;.\;.\;.\;.\;.\;.\;.\;.\;.\;.\;.\;.\;.\;.\;.\;.\;.\;.\; & $271\;6\;7{,}5$\\
\multicolumn{2}{l}{(beide vom mittleren Äquinoktium aus gezählt)}\\
Logarithmus des kleinsten Abstands .\;.\;.\;.\;.\;.\;.\;.\;.\;.\;.\;.\; & $9{,}8122168$\\
Logarithmus der mittleren Bewegung .\;.\;.\;.\;.\;.\;.\;.\;.\;.\;.\;.\; & $0{,}2418031$\\
\end{tabular}\end{center}


\Needspace{14\baselineskip}
\par\indent Diese Elemente lassen für die Breiten zu Anfang Dezember noch eine Abweichung von einer Minute übrig. Durch Verteilung ließe sich dieser Fehler noch beträchtlich vermindern. Herr \textsc{Bessel} zog es jedoch vor, die parabolische Hypothese beiseitezulassen und die Beobachtungen möglichst genau mit einer elliptischen Bahn in Einklang zu bringen. Auch wenn es in der Natur der Sache liegt, dass sich die Abmessungen dieser Bahn nicht genau bestimmen lassen, scheint daraus doch mit ziemlicher Sicherheit hervorzugehen, dass die Bahn tatsächlich elliptisch ist und eine hyperbolische Bahn ausscheidet. Nur bei wenigen Kometen lässt sich nach einer einzigen Erscheinung auch nur so viel mit voller Gewissheit behaupten. Die Elemente der von Herrn \textsc{Bessel} bestimmten elliptischen Bahn lauten:


\begin{center}\fontsize{9}{11}\selectfont\setlength{\tabcolsep}{3pt}\renewcommand{\arraystretch}{1.15}\begin{tabular}{lr}
Zeit des Durchgangs durch das Perihel September .\;.\;.\;.\;.\;.\; & $18{,}74986$\\
Neigung der Bahn .\;.\;.\;.\;.\;.\;.\;.\;.\;.\;.\;.\;.\;.\;.\; & $63^\circ\;10'\;53''2$\\
Aufsteigender Knoten .\;.\;.\;.\;.\;.\;.\;.\;.\;.\;.\;.\;.\;.\; & $266\;46\;3{,}1$\\
Perihel .\;.\;.\;.\;.\;.\;.\;.\;.\;.\;.\;.\;.\;.\;.\;.\;.\;.\; & $271\;56\;0{,}1$\\
\multicolumn{2}{l}{(beide vom mittleren Äquinoktium aus gezählt.)}\\
Logarithmus des kleinsten Abstands .\;.\;.\;.\;.\;.\;.\;.\;.\;.\;.\; & $9{,}8105558$\\
Logarithmus der mittleren Bewegung .\;.\;.\;.\;.\;.\;.\;.\;.\;.\;.\; & $0{,}2442946$\\
Exzentrizität .\;.\;.\;.\;.\;.\;.\;.\;.\;.\;.\;.\;.\;.\;.\;.\;.\;.\;.\; & $0{,}9958626$\\
Große Halbachse .\;.\;.\;.\;.\;.\;.\;.\;.\;.\;.\;.\;.\;.\;.\;.\;.\;.\; & $156{,}253$\\
Umlaufszeit .\;.\;.\;.\;.\;.\;.\;.\;.\;.\;.\;.\;.\;.\;.\;.\;.\;.\;.\;.\; & $1953{,}2\text{ Jahre.}$\\
\end{tabular}\end{center}
